% =====================================================================
%  SN7 (Supplementary Notes): Dimension coverage over the annotated corpus
% =====================================================================

\section{Dimension coverage over the annotated corpus}
\label{sec:supp-coverage-note}

This note reports how often each dimension is populated across the 28
\GeneticEvidence{} items in the pilot corpus and how often each conditional
dimension is applicable under its activation rules. The counts are computed
from the released YAML annotations by
\fpath{scripts/compute_coverage.py}. A per-paper report is available in the
repository at
\fpath{annotations/coverage.md}.\footnote{\url{https://github.com/ForomePlatform/genetic-evidence-model/blob/master/annotations/coverage.md}}

\begin{table}[H]
  \centering
  \small
  \caption[Dimension population across 28 \GeneticEvidence{} items]{Dimension
  population across 28 \GeneticEvidence{} items from six papers
  (Duerr \texttt{v0}, Inouye \texttt{v1}), computed by
  \texttt{compute\_coverage.py}. The left block lists always-required
  dimensions; the right block lists conditional dimensions and the number of
  items in which their activation conditions apply.}
  \label{tab:supp-coverage}
  \begin{tabular}{@{}lcc@{\hspace{1em}}lcc@{}}
    \toprule
    \multicolumn{3}{c}{Always-required dimensions} &
    \multicolumn{3}{c}{Conditional dimensions} \\
    Dimension             & Populated & Applicable &
    Dimension             & Populated & Applicable \\
    \midrule
    Knowledge Domain      & 27        & 28         &
    Variant Ascertainment & 14        & 18         \\
    Method                & 28        & 28         &
    Mode of Inheritance   & 0         & 1          \\
    Target Type           & 28        & 28         &
    Mendelian Segregation & 0         & 1          \\
    Resolution            & 28        & 28         &
    Penetrance            & 2         & 4          \\
    Credibility           & 28        & 28         &
    Measurement Target    & 10        & 11         \\
    Phenotype Scale       & 28        & 28         &
    Gene Relation         & 1         & 11         \\
                          &           &            &
    Organism              & 8         & 8          \\
                          &           &            &
    Knockout Type         & 1         & 6          \\
    \bottomrule
  \end{tabular}
\end{table}

The always-required dimensions are populated in all 28 items by
construction, except Knowledge Domain, which is populated in 27. The
exception is a cited anti-p40 antibody trial in Duerr, which concerns
drug-response evidence with no fitting genetic Knowledge Domain.

Among the conditional dimensions, Variant Ascertainment is populated in 14 of
the 18 items for which its activation condition holds. The four Inouye
\texttt{v1} items use \dimval{Variant} as a placeholder Target Type but mark
Variant Ascertainment as not applicable, because a polygenic score has no
single ascertainment mode. This is the CE-IN1 gap. Well-calibrated conditions
such as Organism, populated in 8 of 8 applicable items, and Measurement
Target, populated in 10 of 11, are generally populated whenever their
conditions hold. In contrast, Gene Relation, populated in 1 of 11 applicable
items, and Knockout Type, populated in 1 of 6, expose enumeration and
over-broad-activation gaps that motivate candidate extensions.

Mode of Inheritance and Mendelian Segregation are each populated in 0 of 1
applicable items. The single \dimval{Human Genetics}~$+$~\dimval{Gene} item,
Davis's case--control burden test, addresses both as not applicable. This
indicates that the current activation condition is broader than the context in
which those dimensions are meaningful.
